% foundations/ellipsoid_geometry.tex

\chapter{Curvature on the Ellipsoid}
\label{chapter:ellipsoid_curvature}

Railway alignment modelling is often carried out in projected coordinate
systems such as Gauss--Krüger or UTM. These projections are convenient for
engineering computation and data exchange, but they introduce geometric
distortions and therefore do not represent the physical geometry of the
railway on the Earth exactly. A more intrinsic formulation describes the
alignment directly on the reference surface, typically modelled as an
ellipsoid.

% ============================================================
\section{Curve on a Surface}
% ============================================================

\begin{definition}[Reference ellipsoid]
Let \(\mathcal{E} \subset \mathbb{R}^3\) denote a reference ellipsoid
representing the Earth.
\end{definition}

\begin{definition}[Alignment on the ellipsoid]
An alignment is defined as a curve
\[
\gamma : [0,L] \to \mathcal{E}
\]
parameterized by arc-length \(s\).
\end{definition}

The parameterization satisfies
\[
\|\gamma'(s)\| = 1,
\]
so that \(s\) represents physical distance along the curve.

% ============================================================
\section{Tangent and Normal Decomposition}
% ============================================================

\begin{definition}[Tangent vector]
The tangent vector is defined as
\[
\mathbf{t}(s)=\gamma'(s).
\]
\end{definition}

\begin{definition}[Surface normal]
Let \(\mathbf{n}_{\mathcal{E}}(s)\) denote the unit normal vector of the
ellipsoid at \(\gamma(s)\).
\end{definition}

The second derivative \(\gamma''(s)\) can be decomposed into a component
tangent to the surface and a component normal to the surface. This gives

\begin{definition}[Decomposition]
\[
\gamma''(s)
=
\underbrace{\nabla_s \mathbf{t}(s)}_{\text{tangential}}
+
\underbrace{(\gamma''(s)\cdot \mathbf{n}_{\mathcal{E}})\,
\mathbf{n}_{\mathcal{E}}}_{\text{normal}}.
\]
\end{definition}

Here, \(\nabla_s\mathbf{t}\) denotes the covariant derivative of the
tangent vector along the curve on the surface.

% ============================================================
\section{Geodesic Curvature}
% ============================================================

\begin{definition}[Geodesic curvature]
The intrinsic curvature of the alignment on the ellipsoid is defined as
\[
\kappa_g(s)
=
\bigl\|\nabla_s \mathbf{t}(s)\bigr\|.
\]
\end{definition}

Geodesic curvature measures the change of direction within the surface.
It is the surface-intrinsic counterpart of planar curvature and is the
relevant curvature notion for alignment geometry constrained to the
reference surface.

% ============================================================
\section{Normal Curvature}
% ============================================================

\begin{definition}[Normal curvature]
The normal curvature is defined as the normal component of the second
derivative:
\[
\kappa_n(s)
=
\bigl|
\gamma''(s)\cdot \mathbf{n}_{\mathcal{E}}(s)
\bigr|.
\]
\end{definition}

Normal curvature is induced by the curvature of the ellipsoid itself and
does not describe the bending of the alignment within the surface.

% ============================================================
\section{Total Curvature}
% ============================================================

\begin{proposition}
The total spatial curvature satisfies:
\[
\kappa^2(s)
=
\kappa_g^2(s)+\kappa_n^2(s).
\]
\end{proposition}

Thus, the spatial curvature decomposes into an intrinsic surface
contribution and an extrinsic surface-induced contribution.

% ============================================================
\section{Railway-Relevant Curvature}
% ============================================================

\begin{proposition}
For railway alignment modelling on a reference surface, the
railway-relevant curvature is the geodesic curvature:
\[
\kappa_{\mathrm{rail}}(s)=\kappa_g(s).
\]
\end{proposition}

Railway vehicles move along the surface, and dynamic quantities such as
lateral acceleration depend on curvature within the surface rather than
on the embedding curvature of the Earth.

% ============================================================
\section{Relation to Planar Models}
% ============================================================

In planar geometry, curvature is defined as
\[
\kappa(s)=\theta'(s).
\]
On a curved surface, the heading-angle derivative is replaced by the
covariant change of the tangent direction:
\[
\kappa_g(s)
=
\bigl\|\nabla_s \mathbf{t}(s)\bigr\|.
\]
Classical planar formulations are therefore a special case of the intrinsic
surface formulation.

% ============================================================
\section{Special Case: Geodesics}
% ============================================================

\begin{definition}[Geodesic]
A curve satisfies
\[
\kappa_g(s)=0
\]
if and only if it is a geodesic on the surface.
\end{definition}

Geodesics represent straightest paths on the ellipsoid in the intrinsic
surface sense. They are generally not straight in projected coordinate
systems. Railway alignment design is therefore not equivalent to geodesic
design; it must also satisfy engineering, dynamic, operational, and
realization constraints.

% ============================================================
\section{Computational Formulation}
% ============================================================

In practical computations, geodesic curvature can be evaluated by
projecting the second derivative onto the tangent plane of the surface:
\[
\kappa_g
=
\left\|\gamma'' - (\gamma''\cdot \mathbf{n}_{\mathcal{E}})
\mathbf{n}_{\mathcal{E}}\right\|.
\]
This formulation separates curvature within the surface from curvature
caused by the embedding of the surface in three-dimensional space.

% ============================================================
\section{Implications for Engineering}
% ============================================================

\begin{proposition}
Curvature evaluated in projected coordinate systems differs from
intrinsic curvature on the ellipsoid.
\end{proposition}

Consequently, arc-length is only approximately preserved, curvature is
distorted by projection, and dynamic quantities are evaluated on an
approximate geometry.

\begin{proposition}
Evaluating curvature intrinsically yields more physically consistent
results.
\end{proposition}

% ============================================================
\section{Vision}
% ============================================================

A consistent modelling framework may formulate alignment geometry
directly on the ellipsoid and use projections only for visualization
and data exchange.

This enables:
\begin{itemize}
\item physically meaningful curvature,
\item accurate dynamic evaluation,
\item improved robustness in optimization,
\item and increased safety in alignment design.
\end{itemize}

This is a long-term research direction rather than a prerequisite for
practical AIM implementations.

% ============================================================
\section{Connection to AIM}
% ============================================================

\begin{summary}
Curvature on the ellipsoid must be understood as an intrinsic quantity,
measured within the surface rather than in ambient space.

For railway alignment on a reference surface, the relevant curvature is
the geodesic curvature.

This distinction is essential for physically consistent modelling and
provides a foundation for extending the AIM framework beyond
projection-based engineering geometry.
\end{summary}
